\section{引言：拆开 Coding Agent 的黑盒}

本期继续探索 Codex，从前几期的 Context Engineering、Runtime、Harness 角度之后，本期聚焦于 Codex 内部\textbf{12 个预定义工具}（Tools）的设计。讲者反复强调的核心理念是：\textbf{没有魔法，千万不要把 Agent 当黑盒}——要去 Review 和 Guide Agent 的内部处理过程，尤其当模型陷入死循环或表现不符合预期时。

\begin{knowledgebox}{Codex 的 12 个预定义 Tool}
\begin{enumerate}
    \item Apply Patch（编辑代码文件）
    \item Execute Commands（执行命令）
    \item Write Stdin（标准输入写入）
    \item Update Plan（更新计划）
    \item Request User Input（请求用户输入）
    \item Web Search（网络搜索）
    \item View Image（查看图像）
    \item Spawn Agents（创建子 Agent）
    \item Send Inputs（向子 Agent 发送指令）
    \item Wait（等待子 Agent）
    \item Close Agents（关闭子 Agent）
    \item Resume（恢复子 Agent）
\end{enumerate}
\end{knowledgebox}


